\section{Introduction}

% 城市客流预测是智能交通系统中的基础任务之一，能够为交通调度、运力配置、拥堵缓解和公共安全管理提供关键支撑。随着城市传感器、交通卡口、共享出行和移动设备数据的持续积累，深度学习方法已经被广泛用于刻画客流序列中的时间周期性、空间相关性和长期动态变化。现有时空预测模型在常规通勤场景下取得了显著进展，但它们大多依赖相对稳定且充足的历史模式，默认未来客流会延续日周期、周周期或局部空间依赖等规律。因此，当城市运行状态偏离常规模式时，仅依靠常规历史模式进行预测往往难以满足实际需求。

Urban flow forecasting is one of the fundamental tasks in intelligent transportation systems, providing essential support for traffic scheduling, capacity allocation, congestion mitigation, and public safety management. With the continuous accumulation of data from urban sensors, traffic checkpoints, shared mobility, and mobile devices, deep learning methods have been widely adopted to characterize the temporal periodicity, spatial correlations, and long-term dynamics of flow sequences. Existing spatiotemporal forecasting models have achieved remarkable progress in regular commuting scenarios, but they largely rely on relatively stable and abundant historical patterns, implicitly assuming that future flows will follow daily/weekly periodicities and local spatial dependencies. Consequently, when the operating state of a city deviates from its regular patterns, forecasts that rely solely on regular historical patterns often fall short of practical requirements.

% 异常场景下的客流预测更加困难，也更具有现实意义。节假日、极端天气、大型活动和突发事件等因素会显著改变居民出行需求，使客流在时间强度、空间分布和变化趋势上都偏离正常通勤模式。针对这类问题，已有研究尝试引入事件信息、天气特征、鲁棒时序建模或异常增强的时空结构，以提升模型对稀有高影响时段的预测能力。这些方法在单城市异常预测中具有一定效果，但通常仍在绝对客流空间中直接建模，并且主要关注训练和测试分布相对一致的场景。由此带来的问题是，模型容易把城市本身的正常客流尺度、局部空间结构和异常造成的偏移混合在一起，难以区分哪些模式是城市特有的，哪些异常变化具有跨城市迁移价值。

Flow forecasting under anomalous scenarios is more difficult and of greater practical importance. Holidays, extreme weather, large-scale events, and emergencies can substantially reshape travel demand, driving flows away from normal commuting patterns in temporal intensity, spatial distribution, and evolving trends. To address such problems, prior studies have incorporated event information, weather features, robust temporal modeling, or anomaly-enhanced spatiotemporal structures to improve forecasting on rare but high-impact periods. These methods are effective for single-city anomaly forecasting to some extent, but they typically still model flows directly in the absolute-value space and mainly target scenarios where the training and test distributions are relatively consistent. As a consequence, models tend to entangle a city's inherent normal flow scale, local spatial structures, and anomaly-induced shifts, making it difficult to distinguish which patterns are city-specific and which anomalous variations are transferable across cities.

% 本文进一步关注一个更具体也更具挑战性的研究问题：跨城市异常客流预测，即模型在源城市上训练，并直接迁移到目标城市，在全量时段和异常时段上进行评估。这个问题具有实际价值，因为单个城市中的异常样本往往数量有限，而不同城市之间可能共享某些由节假日、天气或活动引发的异常变化规律。然而，直接进行跨城市迁移面临两个核心挑战。第一，异常客流样本稀疏，且与正常客流差异显著；模型在大量正常样本上训练时，容易低估异常偏移，或者将异常变化吸收到城市特定的正常基线中。第二，异常客流模式在不同城市之间分布不均，同类异常事件在发生频率、影响强度、空间覆盖范围和持续时间上都可能存在明显差异，因此在原始客流空间中直接迁移往往并不可靠。

This paper further focuses on a more specific and more challenging research problem: cross-city anomalous urban flow forecasting, where a model is trained on a source city and directly transferred to a target city for evaluation on both full and anomalous periods. This problem is of practical value because anomalous samples within a single city are often limited, whereas different cities may share certain anomalous variation patterns induced by holidays, weather, or events. However, direct cross-city transfer faces two core challenges. First, anomalous flow samples are scarce and differ markedly from normal flows; when trained on abundant normal samples, a model tends to underestimate anomalous shifts or absorb anomalous variations into the city-specific normal baseline. Second, anomalous flow patterns are unevenly distributed across cities; the same type of anomalous event may differ substantially in frequency, impact intensity, spatial coverage, and duration, so direct transfer in the raw flow space is often unreliable.

% 我们的核心观察是，城市客流可以被分解为城市特定的正常基线和相对更可迁移的异常偏差。相比直接迁移不同城市之间尺度差异显著的绝对客流值，我们首先利用正常基线刻画目标城市自身的常规出行模式，再在相对于该基线的偏差空间中学习异常动态。基于这一思想，本文提出 **ProtoDRIFT**，一个面向跨城市异常客流预测的原型解耦偏差参照不变预测框架。ProtoDRIFT 包含两个关键设计：其一是可迁移异常原型机制，通过紧凑的原型库归纳偏差空间中反复出现的异常变化模式，从而捕获较少依赖具体城市尺度的异常语义；其二是互信息约束的解耦设计，通过约束异常表示与城市特定表示之间的冗余依赖，促使模型区分可迁移异常偏差和本地城市因素。二者结合后，模型能够以目标城市正常基线为参照，叠加迁移得到并经过城市条件调整的异常偏差，从而完成目标城市客流预测。

Our key observation is that urban flows can be decomposed into a city-specific normal baseline and relatively more transferable anomaly deviations. Instead of directly transferring absolute flow values whose scales differ substantially across cities, we first characterize the target city's own regular travel patterns with a normal baseline, and then learn anomalous dynamics in the deviation space relative to this baseline. Based on this idea, we propose **ProtoDRIFT**, a prototype-disentangled, deviation-referenced invariant forecasting framework for cross-city anomalous urban flow forecasting. ProtoDRIFT features two key designs. The first is a \textbf{transferable anomaly prototype mechanism}, which summarizes recurring anomalous variation patterns in the deviation space with a compact prototype bank, thereby capturing anomaly semantics that are less dependent on city-specific scales. The second is a \textbf{mutual-information-constrained disentanglement design}, which restrains the redundant dependence between anomaly representations and city-specific representations, encouraging the model to distinguish transferable anomaly deviations from local city factors. Combining the two, the model produces target-city flow forecasts by taking the target city's normal baseline as a reference and superimposing the transferred, city-conditionally adjusted anomaly deviations.


% 本文的主要贡献包括三点。第一，我们将跨城市异常客流预测建模为一个偏差参照的迁移问题，提出“城市特定正常基线 + 可迁移异常偏差”的预测视角，以缓解绝对客流空间中城市尺度差异带来的迁移困难。第二，我们设计了 ProtoDRIFT 框架，通过可迁移异常原型和互信息约束的解耦机制，对稀疏异常样本中的共享偏差模式进行结构化建模。第三，我们在 NYC 与 CHI 的双向跨城市迁移设置下进行了系统实验，覆盖 `NYC→CHI` 和 `CHI→NYC` 两个 setting，并同时评估 full 与 anomaly 子集，通过主实验、消融实验、参数敏感性和可视化分析验证所提出设计的有效性。


Our main contributions are threefold.

\begin{itemize}
    \item \textbf{Conceptual insight.} We formulate cross-city anomalous urban flow forecasting as a deviation-referenced transfer problem and propose the forecasting perspective of "city-specific normal baseline + transferable anomaly deviation" to alleviate the transfer difficulty caused by city-scale discrepancies in the absolute flow space.
    \item \textbf{Methodological framework.} We design the ProtoDRIFT framework, which structurally models the shared deviation patterns in sparse anomalous samples via transferable anomaly prototypes and a mutual-information-constrained disentanglement mechanism. 
    \item \textbf{Empirical validation.} we conduct systematic experiments under the bidirectional cross-city transfer setting between NYC and CHI, covering both `NYC$\rightarrow$CHI` and `CHI$\rightarrow$NYC` settings and evaluating both the full and anomaly subsets, and validate the effectiveness of the proposed designs through main results, ablation studies, hyperparameter sensitivity, and visualization analyses.
\end{itemize}